Đây là phần trung tâm của workshop. Bạn sẽ huấn luyện một mô hình gợi ý
và đưa vào phục vụ thời gian thực.

\vspace{0.5em}\noindent\textbf{Hiểu về dữ liệu trước khi bắt
tay}\par

Amazon Personalize học từ \textbf{hành vi}, không phải từ thuộc tính sản
phẩm. Nó cần một bảng ghi lại ai đã tương tác với sản phẩm nào, vào lúc
nào.

Định dạng bắt buộc gồm ba cột, tên viết hoa:

\begin{verbatim}
USER_ID,ITEM_ID,EVENT_TYPE,TIMESTAMP
user-133,prod-080,view,1779328787
user-133,prod-008,view,1779328979
user-133,prod-008,add_to_cart,1779329120
user-133,prod-008,purchase,1779329340
\end{verbatim}

\texttt{TIMESTAMP} phải là số nguyên Unix, không phải chuỗi ngày tháng.

\textbf{Bài học đắt giá nhất của workshop này.} Bộ dữ liệu đầu tiên
chúng tôi dùng gồm 4.491 lượt tương tác sinh \textbf{ngẫu nhiên đều}. Mô
hình huấn luyện xong cho kết quả rất tệ, và khi phân tích lại thì lý do
hiển nhiên: thuật toán lọc cộng tác hoạt động bằng cách tìm các nhóm
người dùng có hành vi giống nhau. Nếu ai cũng tương tác ngẫu nhiên thì
không tồn tại nhóm nào để tìm.

Đổi thuật toán không cứu được điều đó. Phải sửa dữ liệu.

Bộ dữ liệu thứ hai mô phỏng hành vi thật với năm đặc điểm: gom theo
phiên truy cập, có phễu chuyển đổi xem rồi mới thêm giỏ rồi mới mua,
phân phối power-law để một nhóm nhỏ sản phẩm hút phần lớn lượt xem, có
nhịp theo giờ trong ngày, và chia người dùng thành các phân khúc sở
thích khác nhau.

Kết quả trên cùng một recipe:

{\def\LTcaptype{none} % do not increment counter
\begingroup
\small
\setlength{\tabcolsep}{3pt}
\renewcommand{\arraystretch}{1.15}
\begin{longtable}[]{@{}lll@{}}
\toprule\noalign{}
Chỉ số & Dữ liệu ngẫu nhiên & Dữ liệu mô phỏng hành vi \\
\midrule\noalign{}
\endhead
\bottomrule\noalign{}
\endlastfoot
Precision@5 & 0,0889 & 0,4348 \\
NDCG@10 & 0,1799 & 0,6512 \\
MRR@25 & 0,1216 & 0,7130 \\
Coverage & 0,8218 & 0,9505 \\
\end{longtable}
\endgroup
}

\vspace{0.5em}\noindent\textbf{Bước 1. Tạo bucket và tải dữ
liệu}\par

\begin{Shaded}
\begin{Highlighting}[]
\ExtensionTok{aws}\NormalTok{ s3 mb s3://fcj{-}recsys{-}data{-}}\OperatorTok{\textless{}}\NormalTok{tên{-}bạn}\OperatorTok{\textgreater{}}
\ExtensionTok{aws}\NormalTok{ s3 cp interactions.csv s3://fcj{-}recsys{-}data{-}}\OperatorTok{\textless{}}\NormalTok{tên{-}bạn}\OperatorTok{\textgreater{}}\NormalTok{/personalize/}
\ExtensionTok{aws}\NormalTok{ s3 cp items.csv s3://fcj{-}recsys{-}data{-}}\OperatorTok{\textless{}}\NormalTok{tên{-}bạn}\OperatorTok{\textgreater{}}\NormalTok{/personalize/}
\end{Highlighting}
\end{Shaded}

\vspace{0.5em}\noindent\textbf{Bước 2. Cho phép Personalize đọc
bucket}\par

Personalize là một dịch vụ riêng, nó cần được cấp quyền mới đọc được
bucket của bạn.

Tạo role:

\begin{enumerate}
\def\labelenumi{\arabic{enumi}.}
\tightlist
\item
  Vào \href{https://console.aws.amazon.com/iam/home\#/roles}{IAM
  console}, chọn \textbf{Create role}
\item
  \textbf{Trusted entity type}: AWS service. Ở \textbf{Use case}, tìm và
  chọn \textbf{Personalize}
\item
  Gán policy \texttt{AmazonS3ReadOnlyAccess}
\item
  Đặt tên \texttt{PersonalizeS3Role}, chọn \textbf{Create role}
\end{enumerate}

Thêm bucket policy cho phép Personalize đọc. Tạo file
\texttt{bucket-policy.json}:

\begin{Shaded}
\begin{Highlighting}[]
\FunctionTok{\{}
  \DataTypeTok{"Version"}\FunctionTok{:} \StringTok{"2012{-}10{-}17"}\FunctionTok{,}
  \DataTypeTok{"Statement"}\FunctionTok{:} \OtherTok{[}\FunctionTok{\{}
    \DataTypeTok{"Effect"}\FunctionTok{:} \StringTok{"Allow"}\FunctionTok{,}
    \DataTypeTok{"Principal"}\FunctionTok{:} \FunctionTok{\{} \DataTypeTok{"Service"}\FunctionTok{:} \StringTok{"personalize.amazonaws.com"} \FunctionTok{\},}
    \DataTypeTok{"Action"}\FunctionTok{:} \OtherTok{[}\StringTok{"s3:GetObject"}\OtherTok{,} \StringTok{"s3:ListBucket"}\OtherTok{]}\FunctionTok{,}
    \DataTypeTok{"Resource"}\FunctionTok{:} \OtherTok{[}
      \StringTok{"arn:aws:s3:::fcj{-}recsys{-}data{-}\textless{}tên{-}bạn\textgreater{}"}\OtherTok{,}
      \StringTok{"arn:aws:s3:::fcj{-}recsys{-}data{-}\textless{}tên{-}bạn\textgreater{}/*"}
    \OtherTok{]}
  \FunctionTok{\}}\OtherTok{]}
\FunctionTok{\}}
\end{Highlighting}
\end{Shaded}

\begin{Shaded}
\begin{Highlighting}[]
\ExtensionTok{aws}\NormalTok{ s3api put{-}bucket{-}policy }\AttributeTok{{-}{-}bucket}\NormalTok{ fcj{-}recsys{-}data{-}}\OperatorTok{\textless{}}\NormalTok{tên{-}bạn}\OperatorTok{\textgreater{}}\NormalTok{ {-}{-}policy file://bucket{-}policy.json}
\end{Highlighting}
\end{Shaded}

\vspace{0.5em}\noindent\textbf{Bước 3. Tạo dataset
group}\par

\begin{enumerate}
\def\labelenumi{\arabic{enumi}.}
\tightlist
\item
  Mở
  \href{https://ap-southeast-1.console.aws.amazon.com/personalize/home?region=ap-southeast-1}{Amazon
  Personalize console}
\item
  Chọn \textbf{Create dataset group}
\item
  \textbf{Name}: \texttt{fcj-recsys-dsg}, \textbf{Domain}: chọn
  \textbf{Custom}
\item
  Chọn \textbf{Create dataset group and continue}
\end{enumerate}

\begin{figure}
\centering
\pandocbounded{\includegraphics[keepaspectratio,alt={Tạo dataset group}]{5-Workshop/5.6/create-dsg.png}}
\caption{Tạo dataset group}
\end{figure}

\vspace{0.5em}\noindent\textbf{Bước 4. Nạp dữ liệu tương
tác}\par

\begin{enumerate}
\def\labelenumi{\arabic{enumi}.}
\tightlist
\item
  Ở bước \textbf{Create interactions dataset}, chọn \textbf{Create
  dataset}
\item
  \textbf{Dataset name}: \texttt{interactions}
\item
  Chọn \textbf{Create new schema}, đặt tên \texttt{interactions-schema},
  dán nội dung:
\end{enumerate}

\begin{Shaded}
\begin{Highlighting}[]
\FunctionTok{\{}
  \DataTypeTok{"type"}\FunctionTok{:} \StringTok{"record"}\FunctionTok{,}
  \DataTypeTok{"name"}\FunctionTok{:} \StringTok{"Interactions"}\FunctionTok{,}
  \DataTypeTok{"namespace"}\FunctionTok{:} \StringTok{"com.amazonaws.personalize.schema"}\FunctionTok{,}
  \DataTypeTok{"fields"}\FunctionTok{:} \OtherTok{[}
    \FunctionTok{\{} \DataTypeTok{"name"}\FunctionTok{:} \StringTok{"USER\_ID"}\FunctionTok{,} \DataTypeTok{"type"}\FunctionTok{:} \StringTok{"string"} \FunctionTok{\}}\OtherTok{,}
    \FunctionTok{\{} \DataTypeTok{"name"}\FunctionTok{:} \StringTok{"ITEM\_ID"}\FunctionTok{,} \DataTypeTok{"type"}\FunctionTok{:} \StringTok{"string"} \FunctionTok{\}}\OtherTok{,}
    \FunctionTok{\{} \DataTypeTok{"name"}\FunctionTok{:} \StringTok{"EVENT\_TYPE"}\FunctionTok{,} \DataTypeTok{"type"}\FunctionTok{:} \StringTok{"string"} \FunctionTok{\}}\OtherTok{,}
    \FunctionTok{\{} \DataTypeTok{"name"}\FunctionTok{:} \StringTok{"TIMESTAMP"}\FunctionTok{,} \DataTypeTok{"type"}\FunctionTok{:} \StringTok{"long"} \FunctionTok{\}}
  \OtherTok{]}\FunctionTok{,}
  \DataTypeTok{"version"}\FunctionTok{:} \StringTok{"1.0"}
\FunctionTok{\}}
\end{Highlighting}
\end{Shaded}

\begin{enumerate}
\def\labelenumi{\arabic{enumi}.}
\setcounter{enumi}{3}
\tightlist
\item
  Chọn \textbf{Next}, rồi \textbf{Import data directly into datasets}
\item
  Điền:

  \begin{itemize}
  \tightlist
  \item
    \textbf{Import job name}: \texttt{import-interactions}
  \item
    \textbf{IAM role}: chọn \texttt{PersonalizeS3Role}
  \item
    \textbf{Data location}:
    \texttt{s3://fcj-recsys-data-\textless{}tên-bạn\textgreater{}/personalize/interactions.csv}
  \end{itemize}
\item
  Chọn \textbf{Start import}
\end{enumerate}

Quá trình nạp mất khoảng 10 đến 20 phút. Trạng thái phải chuyển sang
\textbf{Active}.

Trong lúc chờ, tranh thủ chụp màn hình các bước đã làm cho phần báo cáo.

\vspace{0.5em}\noindent\textbf{Bước 5. Huấn luyện
solution}\par

\begin{enumerate}
\def\labelenumi{\arabic{enumi}.}
\tightlist
\item
  Ở menu trái, chọn \textbf{Custom resources} rồi \textbf{Solutions and
  recipes}
\item
  Chọn \textbf{Create solution}
\item
  Điền:

  \begin{itemize}
  \tightlist
  \item
    \textbf{Solution name}: \texttt{fcj-user-personalization}
  \item
    \textbf{Solution type}: \textbf{Item recommendation}
  \item
    \textbf{Recipe}: \texttt{aws-user-personalization-v2}
  \end{itemize}
\item
  Chọn \textbf{Create and train solution}
\end{enumerate}

Huấn luyện mất khoảng 40 đến 90 phút tuỳ lượng dữ liệu.

\vspace{0.5em}\noindent\textbf{Bước 6. Xem chỉ số đánh
giá}\par

Khi solution version chuyển sang \textbf{Active}, mở nó ra để xem các
chỉ số:

\begin{itemize}
\tightlist
\item
  \textbf{Precision@K}: tỉ lệ sản phẩm đúng trong K gợi ý đầu tiên
\item
  \textbf{NDCG@K}: giống Precision nhưng có tính vị trí, gợi ý đúng nằm
  đầu danh sách có giá trị cao hơn nằm cuối
\item
  \textbf{MRR@K}: vị trí trung bình của gợi ý đúng đầu tiên
\item
  \textbf{Coverage}: tỉ lệ sản phẩm trong kho từng được gợi ý
\end{itemize}

\begin{figure}
\centering
\pandocbounded{\includegraphics[keepaspectratio,alt={Chỉ số mô hình}]{5-Workshop/5.6/metrics.png}}
\caption{Chỉ số mô hình}
\end{figure}

\vspace{0.5em}\noindent\textbf{Bước 7. Tạo
campaign}\par

Campaign là điểm cuối phục vụ suy luận thời gian thực.

\begin{enumerate}
\def\labelenumi{\arabic{enumi}.}
\tightlist
\item
  Chọn \textbf{Campaigns} rồi \textbf{Create campaign}
\item
  Điền:

  \begin{itemize}
  \tightlist
  \item
    \textbf{Campaign name}: \texttt{fcj-recsys-campaign}
  \item
    \textbf{Solution}: chọn solution vừa huấn luyện
  \item
    \textbf{Minimum provisioned TPS}: \texttt{1}
  \end{itemize}
\item
  Chọn \textbf{Create campaign}
\end{enumerate}

Campaign bắt đầu \textbf{tính phí theo giờ} ngay khi tạo, kể cả khi
không có truy vấn nào. Đây là tài nguyên tốn tiền nhất trong workshop.
Ghi nhớ để xoá ở phần dọn dẹp.

\vspace{0.5em}\noindent\textbf{Bước 8. Nối campaign vào
Lambda}\par

Copy Campaign ARN, rồi thêm biến môi trường cho Lambda:

\begin{enumerate}
\def\labelenumi{\arabic{enumi}.}
\tightlist
\item
  Mở hàm \texttt{fcj-recsys-api}, tab \textbf{Configuration}, mục
  \textbf{Environment variables}
\item
  Chọn \textbf{Edit}, thêm biến:

  \begin{itemize}
  \tightlist
  \item
    \textbf{Key}: \texttt{CAMPAIGN\_ARN}
  \item
    \textbf{Value}: ARN vừa copy
  \end{itemize}
\item
  Chọn \textbf{Save}
\end{enumerate}

\vspace{0.5em}\noindent\textbf{Kiểm tra}\par

\begin{Shaded}
\begin{Highlighting}[]
\ExtensionTok{aws}\NormalTok{ personalize{-}runtime get{-}recommendations }\DataTypeTok{\textbackslash{}}
  \AttributeTok{{-}{-}campaign{-}arn} \OperatorTok{\textless{}}\NormalTok{CAMPAIGN{-}ARN}\OperatorTok{\textgreater{}} \DataTypeTok{\textbackslash{}}
  \AttributeTok{{-}{-}user{-}id}\NormalTok{ user{-}001 }\DataTypeTok{\textbackslash{}}
  \AttributeTok{{-}{-}num{-}results}\NormalTok{ 5}
\end{Highlighting}
\end{Shaded}

Trả về danh sách \texttt{itemId} là mô hình đã hoạt động. Thử với vài
\texttt{user-id} khác nhau, kết quả phải khác nhau --- đó là bằng chứng
gợi ý thực sự được cá nhân hoá.
